Although SPSS Modeler provides a structured visual environment that aligns well with many CRISP-DM activities, several practical limitations were encountered during this project. In particular, not every CRISP-DM phase can be represented directly as a modelling node. Activities such as business understanding, interpretation of clinical requirements, and deployment planning still require substantial documentation and reasoning outside the
SPSS stream.

Model refinement was also relatively manual. Although SPSS provides automated modelling facilities for comparing multiple algorithms, detailed hyperparameter tuning often requires individual model nodes to be configured,
executed, and compared repeatedly. This makes systematic tuning less flexible than programmatic frameworks where parameter searches, cross-validation, and experiment tracking can be automated more easily.

The range of algorithms and transformations available also depends on the
installed SPSS Modeler version and extensions. Some newer machine-learning methods and specialised preprocessing techniques are less readily available
than in Python or R ecosystems. Similarly, custom feature engineering and complex pipeline logic can become difficult to maintain when a stream contains
many connected nodes.

Another challenge concerns reproducibility and workflow management. The visual stream is useful for demonstrating the sequence of CRISP-DM activities, but changes to parameters and intermediate nodes must be carefully documented to
ensure that modelling iterations can be reproduced. For this reason, key settings, transformations, partitions, selected variables, and evaluation
results were explicitly recorded in the report.

Finally, deployment in this iteration is limited to demonstrating a model-ready SPSS workflow rather than implementing a production clinical
application. A real deployment would additionally require automated data ingestion, monitoring, model versioning, security controls, and integration with healthcare information systems.

Despite these limitations, SPSS Modeler was appropriate for this iteration because its visual workflow clearly demonstrates the relationship between data understanding, preparation, transformation, modelling, and evaluation.
The limitations mainly affect scalability, automation, and production deployment rather than the ability to conduct the current proof-of-concept
classification analysis.